\documentclass[a4paper,11pt]{article}

\usepackage[T1]{fontenc}
\usepackage[utf8]{inputenc}
\usepackage{lmodern}
\usepackage{microtype}
\usepackage{geometry}
\usepackage{fancyhdr}
\usepackage{lastpage}
\usepackage{amsmath,amssymb}
\usepackage{booktabs}
\usepackage{array}
\usepackage{tabularx}
\usepackage{pdflscape}
\usepackage{enumitem}
\usepackage{xcolor}
\usepackage{listings}
\usepackage[hidelinks]{hyperref}

\geometry{margin=2.5cm}
\setlength{\headheight}{13.6pt}
\setlength{\parindent}{0pt}
\setlength{\parskip}{0.6em}

\newcommand{\coursename}{SH2404 Astrophysics}
\newcommand{\reporttitle}{Observations and Position Calculations}
\newcommand{\studentname}{Lorenzo Deflorian}
\newcommand{\observationdate}{1 and 10 May 2026}
\newcommand{\observationtime}{21:30 and 23:30 local time}
\newcommand{\calculationdate}{14 February 2026}
\newcommand{\calculationtime}{18:00 local time}
\newcommand{\observerlatitude}{$59.3900051^\circ$ N}
\newcommand{\observerlongitude}{$18.006606^\circ$ E}

\definecolor{codebackground}{RGB}{248,248,248}
\definecolor{codeborder}{RGB}{210,210,210}
\definecolor{codekeyword}{RGB}{0,70,140}
\definecolor{codecomment}{RGB}{90,110,90}

\lstdefinestyle{reportpython}{
    language=Python,
    basicstyle=\ttfamily\small,
    keywordstyle=\color{codekeyword}\bfseries,
    commentstyle=\color{codecomment}\itshape,
    backgroundcolor=\color{codebackground},
    frame=single,
    rulecolor=\color{codeborder},
    framerule=0.4pt,
    breaklines=true,
    columns=fullflexible,
    keepspaces=true,
    showstringspaces=false,
    tabsize=4,
    literate={°}{{$^\circ$}}1,
    xleftmargin=0.8em,
    xrightmargin=0.8em,
    framexleftmargin=0.8em,
    framexrightmargin=0.8em
}

\newcolumntype{L}[1]{>{\raggedright\arraybackslash}p{#1}}
\newcolumntype{Y}{>{\raggedright\arraybackslash}X}

\pagestyle{fancy}
\fancyhf{}
\lhead{\coursename}
\rhead{\thepage\ of \pageref{LastPage}}
\renewcommand{\headrulewidth}{0.4pt}

\begin{document}
\pagenumbering{roman}

\begin{titlepage}
    \centering
    \vspace*{2cm}

    {\Large \coursename\par}
    \vspace{1.5cm}

    {\huge\bfseries \reporttitle\par}
    \vspace{2cm}

    \begin{tabular}{rl}
        Student: & \studentname \\
        Observation date: & \observationdate \\
        Observation time: & \observationtime \\
        Observer latitude: & \observerlatitude \\
        Observer longitude: & \observerlongitude \\
        Report date: & \today \\
    \end{tabular}

    \vfill
\end{titlepage}

\pagenumbering{arabic}
\tableofcontents
\clearpage

\section{Abstract}
This report combines naked-eye observations from Ulriksdal Palace park with simple position calculations for selected astronomical objects. The observation logbook covers two observing times on 1 May 2026 and a comparison observation on 10 May 2026. The clearest observed targets were the Moon, Jupiter, Ursa Major, and Cassiopeia, while faint objects such as the Milky Way and the Pleiades were not visible because of twilight, moonlight, light pollution, and a partly obstructed horizon. The calculation part uses sidereal time, hour angle, and declination limits to estimate where selected objects should be in the sky from Stockholm. The calculations and Stellarium check show that Sirius and Jupiter are well placed in the evening calculation case, while Mimosa never rises from Stockholm and several other objects depend strongly on time, horizon obstruction, and sky brightness.

\section{Introduction}
The purpose of this exercise is to connect direct naked-eye observations of the night sky with simple calculations of object positions. The first part of the work consists of observing the sky on different nights and recording visible constellations, planets, the Moon, weather conditions, and changes over time. The second part estimates which selected astronomical objects should be visible from Stockholm by using right ascension, declination, sidereal time, hour angle, and altitude limits.

My own observations were made from the park near Ulriksdal Palace, behind the gardens, at approximately $59.3900051^\circ$ N and $18.006606^\circ$ E. I chose this place because the trees reduced some direct light pollution, while the site was still accessible in the evening. The same trees also blocked the lowest part of the horizon, so objects very close to the horizon were difficult to judge.

\section{Observation Logbook}
\label{sec:logbook}

The three naked-eye observing sessions are written below as separate notes. The two observations on 1 May were separated by two hours, which made it possible to compare the apparent motion of the sky during a single night. The observation on 10 May was used as a later comparison under different Moon conditions.

\subsection*{Observation 1: 1 May 2026, 21:30}
\textbf{Site:} Ulriksdal Palace park, $59.3900051^\circ$ N, $18.006606^\circ$ E.

\textbf{Weather and sky conditions:} The sky was clear, but bright twilight affected the western and north-western sky. The Moon was bright, the sky background was therefore bright, and the lowest part of the horizon was partly blocked by trees.

\textbf{Equipment:} Naked-eye observation.

\textbf{Objects and notes:} The Big Dipper, in Ursa Major, was high toward the north-west/west and easy to identify. Cassiopeia was lower toward the north-east, with its W shape visible but less clear. The Moon was low toward the south-east/east and strongly reduced the contrast of the surrounding sky. Jupiter was visible as a bright object in the west/north-west and became easier to see as the sky darkened. Orion, Sirius, the Pleiades, the Milky Way, Venus, and Mars were not confidently visible.

\subsection*{Observation 2: 1 May 2026, 23:30}
\textbf{Site:} Ulriksdal Palace park, same site.

\textbf{Weather and sky conditions:} The sky was still clear enough for observation and darker than at 21:30, but the bright Moon and local light pollution still limited faint-object visibility. Trees continued to block the lowest part of the horizon.

\textbf{Equipment:} Naked-eye observation.

\textbf{Objects and notes:} The Big Dipper was still high and clear, but it had shifted westward. Cassiopeia was visible toward the north-east/north and had also shifted slightly. The Moon was higher than before and closer to the southern sky. Jupiter remained visible, but it was lower toward the west/north-west. Orion, Sirius, the Pleiades, the Milky Way, Venus, and Mars were not seen.

\subsection*{Observation 3: 10 May 2026, 21:30}
\textbf{Site:} Ulriksdal Palace park, same site.

\textbf{Weather and sky conditions:} The early-evening sky was usable, with no major clouds blocking the northern sky. Strong late-spring twilight and local light pollution were still present, but the Moon was not visible during the observation, so the sky was darker than on 1 May.

\textbf{Equipment:} Naked-eye observation.

\textbf{Objects and notes:} The Big Dipper was the clearest constellation and was high in the sky. Cassiopeia was lower toward the north-east and less prominent than Ursa Major. Jupiter was visible after sunset and was one of the easiest bright objects to identify. Orion, Sirius, the Pleiades, the Milky Way, Venus, Mars, and the Moon were not seen.

\subsection{Identified Constellations and Objects}
The Big Dipper, in Ursa Major, was the most reliable constellation marker in all observations. On 1 May at 21:30 it was high toward the north-west/west, and at 23:30 it was still high but had shifted westward. On 10 May it was again high in the sky and was the clearest constellation pattern.

Cassiopeia was visible in all observations, but it was less prominent than the Big Dipper. It was seen lower toward the north-east on 1 May at 21:30, toward the north-east/north at 23:30, and lower toward the north-east on 10 May. The W shape could be recognized, but it was affected by the bright northern sky.

Orion, Sirius, the Pleiades, and the Milky Way were not visible in these observations. Orion and Sirius were no longer well placed in the late spring evening sky, while the Pleiades and the Milky Way were too faint for the sky conditions. Twilight, Stockholm light pollution, moonlight on 1 May, and the partly obstructed horizon all made faint-object observing difficult.

\subsection{Planets and Moon}
Jupiter was the only planet that I could identify confidently. On 1 May at 21:30 it appeared as a bright object in the western/north-western sky and became easier to see as the sky darkened. At 23:30 it was still visible but lower toward the west/north-west. On 10 May it was again visible after sunset, although it was lower as the evening continued.

Venus could not be identified confidently. It was expected to be low in the western sky, but the western horizon was bright and partly obstructed. Mars was not visible during any of the observations.

The Moon was visible on 1 May and was the brightest natural source of sky brightness. At 21:30 it was low toward the south-east/east, and by 23:30 it had moved higher and closer to the southern sky. Its brightness reduced the contrast of the surrounding sky and made faint objects difficult to observe. On 10 May I did not see the Moon during the evening observation, and the sky was therefore darker than on 1 May.

\subsection{Changes Between Observations}
The main change during the two hours on 1 May was the apparent westward motion of the sky caused by Earth's rotation. The Big Dipper shifted westward while remaining high and easy to identify. Cassiopeia also shifted slightly, the Moon moved higher and closer to the southern sky, and Jupiter became lower toward the west/north-west.

The comparison with 10 May showed a different sky background. The Moon was not visible during the 10 May evening observation, so the sky was darker than on 1 May. However, the late-spring twilight and local light pollution still prevented faint targets such as the Milky Way and Pleiades from being seen. The same northern constellations remained the most useful reference patterns.

\section{Preparatory Questions}
\label{sec:questions}

\begin{enumerate}[label=\textbf{\arabic*.}, leftmargin=*]
    \item If a star has an hour angle of 6 h, it is found in the western sky. It is exactly due west on the horizon only if the star lies on the celestial equator.

    \item When a star is on the observer's meridian, the local sidereal time is equal to the star's right ascension:
    \[
        \Theta = \alpha.
    \]

    \item The hour angle of the vernal equinox is the local sidereal time.

    \item The time since an object passed the local meridian is called the hour angle. If the hour angle is zero, the object is on the local meridian; if it is positive, the object has already passed it.

    \item The sidereal time is $0\mathrm{h}\,0\mathrm{m}\,0\mathrm{s}$ when the vernal equinox is on the observer's meridian.

    \item At noon, the Sun is on the observer's local meridian.

    \item Polaris lies about $0.7^\circ$ from the north celestial pole, so its declination is approximately $+89.3^\circ$.

    \item Stockholm is east of the Swedish time meridian, so stars reach the southern meridian earlier in Stockholm than on the time meridian.

    \item A star reaches its highest point when it crosses the local meridian. For a star with declination $30^\circ$, its maximum altitude from Stockholm is approximately
    \[
        h_{\max} = 30^\circ + 90^\circ - 59^\circ 21' = 60^\circ 39'.
    \]
    It is then found due south.

    \item Since a star reaches its highest point due south, its lowest point is found due north.
\end{enumerate}

\section{Calculation Method}
\label{sec:method}

\subsection{Input Data and Assumptions}
The calculation example is planned for \calculationdate\ at \calculationtime. The observer is assumed to be near Stockholm, at latitude \observerlatitude\ and longitude \observerlongitude. The object catalogue used for the computations is shown in Table~\ref{tab:input-objects}.

\begin{table}[htbp]
    \centering
    \caption{Objects used in the observation-planning calculation.}
    \label{tab:input-objects}
    \begin{tabular}{llll}
        \toprule
        Object & Type & Right ascension $\alpha$ & Declination $\delta$ \\
        \midrule
        Sirius & Star & 6h 46m 16s & $-16^\circ 45'$ \\
        Mimosa & Star & 12h 49m 14s & $-59^\circ 50'$ \\
        Albireo & Double star & 19h 31m 43s & $+28^\circ 00'$ \\
        Ring nebula & Nebula & 18h 54m 30s & $+33^\circ 03'$ \\
        Andromeda & Galaxy & 0h 44m 05s & $+41^\circ 24'$ \\
        The Moon & Moon & 19h 40m 42s & $-25^\circ 47'$ \\
        Venus & Planet & 22h 30m 52s & $-10^\circ 54'$ \\
        Jupiter & Planet & 7h 09m 48s & $+22^\circ 47'$ \\
        Mars & Planet & 21h 21m 06s & $+16^\circ 37'$ \\
        \bottomrule
    \end{tabular}
\end{table}

The right ascension and declination values are taken from the lab handout table. The calculations use the simplified approximation from the assignment rather than a full precise ephemeris, so the results should be treated as planning estimates.

\subsection{Sidereal Time and Hour Angle}
\label{sec:computations}
The simplified computation script gives the local sidereal time for the observation as
\[
    \Theta = 04\mathrm{h}\,05\mathrm{m}\,00\mathrm{s}.
\]
For each object, the hour angle is computed from
\[
    H = \Theta - \alpha,
\]
where $\Theta$ is the sidereal time and $\alpha$ is the object's right ascension. The result is wrapped into the range from 0 h to 24 h.

The position estimate uses the following approximate direction mapping:
\[
    H = 0\mathrm{h} \rightarrow \mathrm{south}, \quad
    H = 6\mathrm{h} \rightarrow \mathrm{west}, \quad
    H = 12\mathrm{h} \rightarrow \mathrm{north}, \quad
    H = 18\mathrm{h} \rightarrow \mathrm{east}.
\]

For simplicity, the directions between the four exact hour angles are treated as broad compass sectors rather than as exact azimuths. This is kept like this in the code because the purpose of the calculation is only to estimate where to look in the sky, not to compute a precise telescope pointing.

\begin{lstlisting}[style=reportpython]
def get_direction(sidereal_time: SiderealTime,
                  right_ascension: RightAscension) -> Direction:
    hour_angle = get_hour_angle(sidereal_time, right_ascension)
    # 0h=South, 6h=West, 12h=North and 18h=East
    if hour_angle.total_seconds == 0:
        return Direction.SOUTH
    if hour_angle.total_seconds < 6 * 60 * 60:
        return Direction.SOUTHWEST
    if hour_angle.total_seconds == 6 * 60 * 60:
        return Direction.WEST
    if hour_angle.total_seconds < 12 * 60 * 60:
        return Direction.NORTHWEST
    if hour_angle.total_seconds == 12 * 60 * 60:
        return Direction.NORTH
    if hour_angle.total_seconds < 18 * 60 * 60:
        return Direction.NORTHEAST
    if hour_angle.total_seconds == 18 * 60 * 60:
        return Direction.EAST
    return Direction.SOUTHEAST
\end{lstlisting}

\subsection{Maximum Altitude}
\textbf{Task i.} Show that the maximum height of a star above the horizon is
\[
    h_{\max} = \delta_* + 90^\circ - \mathrm{lat}.
\]

At the highest point of the star, it crosses the local meridian.

The zenith is at latitude $\mathrm{lat}$ above the celestial equator, while the star is at declination $\delta_*$.  
So the angular distance between the zenith and the star is

\[
z = \mathrm{lat} - \delta_*
\]

where $z$ is the zenith distance.

Since altitude and zenith distance add up to $90^\circ$,

\[
h_* + z = 90^\circ
\]

we get

\[
h_* + (\mathrm{lat} - \delta_*) = 90^\circ
\]

and therefore

\[
h_* = \delta_* + 90^\circ - \mathrm{lat}.
\]

\subsection{Minimum Altitude}
\textbf{Task ii.} Use the maximum-height relation to find a relation between the declination of a star and its minimum height above or below the horizon.

At the lowest point of its path, the star is on the opposite side of the celestial pole, toward the north.

The north celestial pole is at an altitude equal to the observer's latitude, \(\mathrm{lat}\). A star with declination \(\delta_*\) is \(90^\circ - \delta_*\) away from the north celestial pole. Therefore its lowest altitude is

\[
h_{\min} = \mathrm{lat} - (90^\circ - \delta_*)
\]

and therefore

\[
h_{\min} = \mathrm{lat} + \delta_* - 90^\circ
\]

The star is circumpolar when

\[
\mathrm{lat} + \delta_* > 90^\circ
\]

In that case its lowest point is still above the horizon. If this value is negative, the star is below the horizon at its lowest point.

\subsection{Altitude Limits and Positions}
\textbf{Task iii.} Using the relations above, the highest and lowest altitude of each object can be computed from
\[
    h_{\max} = \delta_* + 90^\circ - \mathrm{lat},
    \qquad
    h_{\min} = \mathrm{lat} + \delta_* - 90^\circ.
\]

\textbf{Task iv.} The hour angle then gives an estimate of where each object is in the sky at 18:00. Objects near the southern half of the sky are more likely to be visible at that time, while objects on the northern half may be below the horizon unless they are circumpolar.

\begingroup
\setlength{\tabcolsep}{2pt}
\renewcommand{\arraystretch}{1.15}
\begin{table}[htbp]
    \centering
    \scriptsize
    \caption{Computed positions for \calculationdate\ at \calculationtime.}
    \label{tab:computed-objects}
    \begin{tabularx}{\linewidth}{L{1.7cm}L{1.5cm}L{2.0cm}L{1.55cm}L{2.0cm}L{1.35cm}L{1.35cm}L{1.7cm}}
    \toprule
    Object & Type & $\alpha$ & $\delta$ & $H$ & $h_{\max}$ & $h_{\min}$ & Direction \\
    \midrule
    Sirius & Star & 06h 46m 16s & $-16^\circ 45'$ & 21h 18m 44s & 13.9$^\circ$ & -47.4$^\circ$ & Southeast \\
    Mimosa & Star & 12h 49m 14s & $-59^\circ 50'$ & 15h 15m 46s & -29.2$^\circ$ & -90.5$^\circ$ & Northeast \\
    Albireo & Double star & 19h 31m 43s & $+28^\circ 00'$ & 08h 33m 17s & 58.6$^\circ$ & -2.7$^\circ$ & Northwest \\
    Ring nebula & Nebula & 18h 54m 30s & $+33^\circ 03'$ & 09h 10m 30s & 63.7$^\circ$ & 2.4$^\circ$ & Northwest \\
    Andromeda & Galaxy & 00h 44m 05s & $+41^\circ 24'$ & 03h 20m 55s & 72.1$^\circ$ & 10.8$^\circ$ & Southwest \\
    The Moon & Moon & 19h 40m 42s & $-25^\circ 47'$ & 08h 24m 18s & 4.9$^\circ$ & -56.4$^\circ$ & Northwest \\
    Venus & Planet & 22h 30m 52s & $-10^\circ 54'$ & 05h 34m 08s & 19.7$^\circ$ & -41.5$^\circ$ & Southwest \\
    Jupiter & Planet & 07h 09m 48s & $+22^\circ 47'$ & 20h 55m 12s & 53.4$^\circ$ & -7.9$^\circ$ & Southeast \\
    Mars & Planet & 21h 21m 06s & $+16^\circ 37'$ & 06h 43m 54s & 47.3$^\circ$ & -14.0$^\circ$ & Northwest \\
    \bottomrule
    \end{tabularx}
\end{table}
\endgroup

\begin{table}[htbp]
    \centering
    \small
    \caption{Visibility estimates for \calculationdate\ at \calculationtime.}
    \label{tab:visibility-estimates}
    \begin{tabularx}{\linewidth}{L{2.6cm}Y}
    \toprule
    Object & Visibility estimate \\
    \midrule
    Sirius & Probably visible at this time because it is on the southern half of the sky. \\
    Mimosa & Never rises above the horizon from this latitude. \\
    Albireo & Probably not visible at this time, but visible at another time when it is farther south. \\
    Ring nebula & Visible all night because it is circumpolar. \\
    Andromeda & Visible all night because it is circumpolar. \\
    The Moon & Probably not visible at this time, but visible at another time when it is farther south. \\
    Venus & Probably visible at this time because it is on the southern half of the sky. \\
    Jupiter & Probably visible at this time because it is on the southern half of the sky. \\
    Mars & Probably not visible at this time, but visible at another time when it is farther south. \\
    \bottomrule
    \end{tabularx}
\end{table}

The values in Tables~\ref{tab:computed-objects} and~\ref{tab:visibility-estimates} were produced with the computation script included in the appendix. The script represents right ascension and sidereal time in hours, minutes, and seconds, computes $H = \Theta - \alpha$, maps hour angle to an approximate compass direction, and evaluates $h_{\max}$ and $h_{\min}$ from the object's declination and the Stockholm latitude.

\section{Visibility Discussion}
The calculation for \calculationdate\ at \calculationtime\ suggests that Sirius and Jupiter should be on the southern half of the sky and therefore realistically observable if the horizon is clear enough. Sirius is always low from Stockholm, with a maximum altitude of only about $14^\circ$, so it is sensitive to trees, buildings, haze, and twilight. Jupiter is much easier because it is bright and reaches a much larger maximum altitude.

Andromeda and the Ring Nebula are circumpolar from Stockholm according to the altitude limits, so they do not set below the horizon. In practice, Andromeda is a much better naked-eye candidate than the Ring Nebula, because the Ring Nebula is small and faint and normally requires optical aid. Albireo, Mars, and the Moon are not well placed at the calculation time but can be visible at other times when they are closer to the southern part of the sky. Mimosa has a negative maximum altitude from Stockholm and therefore never rises above the horizon.

The real observations show how strongly visibility depends on observing conditions. Even when an object is geometrically above the horizon, it may not be visible if it is low, close to twilight, hidden by trees, or too faint for a light-polluted sky. The May observations also took place much later in the year than the February calculation case, so the visible constellations differed: Ursa Major and Cassiopeia were easy to identify, while Orion and Sirius were not seen.

\section{Stellarium Check}
\label{sec:stellarium}

As a check on the approximate calculations, I used Stellarium with the observing location set to Stockholm and with both the atmosphere and landscape enabled. These settings make the check more realistic than the simple calculation, since objects very close to the horizon or close to twilight may be difficult or impossible to see even if their computed position is formally above the horizon.

\begin{enumerate}[label=\textbf{\arabic*.}, leftmargin=*]
    \item Which stars and planets, if any, are visible?

    At 18:00 on 14 February 2026, Stellarium showed Sirius as visible very low above the south-eastern horizon. Jupiter was also visible in the south-east at a medium altitude. Andromeda was visible toward the west at medium altitude, while the Ring Nebula and Albireo were visible toward the north-west, both low in the sky. Orion was visible toward the south-east at medium altitude, Ursa Major was high in the north-east, and Cassiopeia was very high toward the south-west.

    Venus, Mars, and the Moon were below the horizon in the west and were therefore not visible at 18:00. This agrees well with the calculation for Jupiter, Sirius, Andromeda, and the Ring Nebula. The main difference is Venus: the simplified calculation suggested that Venus might be visible in the south-west, but Stellarium showed it just below the western horizon. This difference is reasonable because the calculation only gives a rough sky sector and altitude range, while Stellarium uses the actual ephemeris and horizon geometry.

    \item Which time does the Sun set?

    Stellarium gave the sunset time in Stockholm on 14 February 2026 as approximately 16:40 local time.

    \item At what time will Sirius be highest in the sky?

    Sirius was highest in the sky at about 21:00 local time. At that time Stellarium gave an azimuth and altitude of approximately
    \[
        \mathrm{Az} = 181^\circ 15'07.1'', \qquad
        \mathrm{Alt} = +13^\circ 56'38.0''.
    \]
    This places Sirius almost exactly due south and low above the horizon, as expected for a star with negative declination observed from Stockholm.

    \item One hour after sunset, which stars and planets are visible?

    One hour after sunset, at about 17:40, Sirius was just above the horizon and away from the strongest twilight. Jupiter was at medium altitude and away from twilight, so it should be visible. Venus was close to the twilight region and near the horizon, making it unlikely to be visible in practice. The Moon and Mars were below the horizon at this time.

    \item When changing the day one day at a time until one month has passed, what do you observe?

    Advancing the date one day at a time while keeping the same clock time showed that the stars and constellations shift westward from night to night. Winter constellations such as Orion and Sirius moved closer to setting at the same clock time, while objects that were initially farther east became better placed. The planets did not keep exactly the same pattern as the stars, because their positions change relative to the background constellations.

    \item If observing Venus on 1 April, what time is best?

    On 1 April 2026, Venus was highest during the daytime, so the best practical observing time was shortly after sunset before it set. Stellarium indicated that a good time was about 21:15 local time, when Venus was toward the west at an altitude of approximately $+02^\circ 22'08.5''$. This is still very low, so the observation would require a clear western horizon.
\end{enumerate}

\section{Reflection on Own Observations}
\label{sec:reflection}

The observations made clear that a formal position above the horizon is not enough for a successful naked-eye observation. The calculation and Stellarium sections were mainly for the 14 February planning case, where Sirius, Jupiter, Andromeda, and some northern objects were expected to be visible. My actual observations were in May, so the sky had changed significantly. Orion and Sirius were no longer visible in the evening observations, while Ursa Major and Cassiopeia were the easiest constellations to use as references.

The two observations on 1 May showed the effect of Earth's rotation directly. Over two hours, Jupiter moved lower toward the west/north-west, the Moon moved higher and toward the southern sky, and the Big Dipper shifted westward while staying high. This matched the expected apparent westward motion of the sky. The comparison with 10 May showed that conditions can matter as much as geometry: without the visible Moon, the sky was darker, but late-spring twilight and city light pollution still made faint objects invisible.

The largest practical limitations were the bright sky background and the partly obstructed horizon. Trees helped block some local lights but also made it hard to judge or see objects close to the horizon, especially toward the west. This affected the attempts to identify Venus and may also have contributed to the difficulty of seeing low objects. Estimating direction and altitude by eye was also approximate, so the logbook positions should be treated as broad sky sectors rather than precise measurements.

\section{Conclusion}
The simplified calculations for Stockholm show that Jupiter and Sirius should be observable in the 14 February evening case, although Sirius remains very low. Andromeda and the Ring Nebula are circumpolar, while Albireo, Mars, and the Moon depend on observing time, and Mimosa is never visible from Stockholm. The May logbook observations showed that the easiest naked-eye targets were Jupiter, the Moon, Ursa Major, and Cassiopeia. Fainter targets, including the Milky Way and the Pleiades, were not visible because of twilight, moonlight, light pollution, and the partly blocked horizon.

\appendix
\section{Example Manual Calculation}
For Sirius, the calculation uses
\[
    \Theta = 04\mathrm{h}\,05\mathrm{m}\,00\mathrm{s},
    \qquad
    \alpha = 06\mathrm{h}\,46\mathrm{m}\,16\mathrm{s}.
\]
The hour angle is
\[
    H = \Theta - \alpha
      = 04\mathrm{h}\,05\mathrm{m}\,00\mathrm{s}
        - 06\mathrm{h}\,46\mathrm{m}\,16\mathrm{s}
      = -02\mathrm{h}\,41\mathrm{m}\,16\mathrm{s}.
\]
Wrapping this into the range 0--24 h gives
\[
    H = 21\mathrm{h}\,18\mathrm{m}\,44\mathrm{s}.
\]
In the simplified direction mapping this places Sirius in the south-eastern sky. With $\delta=-16^\circ45'$ and $\mathrm{lat}=59^\circ21'$, the altitude limits are approximately
\[
    h_{\max}=13.9^\circ,
    \qquad
    h_{\min}=-47.4^\circ.
\]
Sirius can therefore rise above the horizon from Stockholm, but it always stays low and is sensitive to horizon obstruction and atmospheric conditions.

\section{Computation Script}
The numerical values in the calculation tables were produced with the Python script below. The script defines the object catalogue, computes the sidereal time and hour angle, maps the hour angle to a broad compass direction, and calculates the maximum and minimum altitude limits from the object's declination and the observer latitude.

\lstinputlisting[
    style=reportpython,
    caption={Full source code used for the computations},
    label={lst:computations-py}
]{../computations.py}

\end{document}
